\documentclass[12pt,a4paper]{article}
\usepackage[utf8]{inputenc}
\usepackage[T1]{fontenc}
\usepackage[ngerman]{babel}
\usepackage[margin=2.5cm]{geometry}
\usepackage{url}

\setlength{\parindent}{0pt}
\setlength{\parskip}{6pt}

\newcommand{\frage}[1]{\par\vspace{0.7em}\noindent\textbf{#1}\par}
\newcommand{\antwort}[1]{\par\noindent\textit{Antwort:} #1\par}

\begin{document}

\begin{center}
  {\LARGE\bfseries Leitfaden mit Antworten}\\[4pt]
  {\large Markus 15,20--41 · Hauskreis 08.10.2026 · nur für die Leitung}\\[4pt]
  {\footnotesize Die Nummern entsprechen den Fragen im Hauptdokument.}
\end{center}

\subsection*{Sprache}

\frage{1. Welche Begriffe sind im Text unklar oder fremd? Sammelt sie und klärt sie mit Hilfe eines Bibellexikons (z.\,B. Bibel-Lexikon, apps.bibelwissenschaft.de).}
\antwort{Mögliche Sammlung mit Kurz-Erklärung: \textbf{Purpurmantel} -- Soldatenkleidungsstück kaiserlicher Farbe, Spott-Accessoire (\glqq Kaiserkleid\grqq{}); \textbf{Myrrhenwein} -- betäubender Schmerztrunk (vgl. Spr 31,6); \textbf{Golgatha} -- aramäisch \glqq Schädel(stätte)\grqq{}, vermutlich nach der Hügelform benannt; \textbf{dritte/sechste/neunte Stunde} -- jüdische Zeitrechnung ab 6 Uhr morgens, also 9/12/15 Uhr; \textbf{Eli, Eli, lama asabtani} -- aramäisch, Markus übersetzt es im Text; \textbf{Elia} -- Prophet, von dem man am Ende der Zeiten seine Wiederkunft erwartete (Maleachi 3,23); \textbf{Vorhang} -- großer Vorhang vor dem Allerheiligsten im Tempel, das Symbol der Trennung zwischen Gott und Mensch. Wichtig fürs Gespräch: Fragen zu Wortbedeutungen zuerst von der Gruppe sammeln lassen -- sie klären viele Verständnisprobleme, bevor die eigentliche Analyse beginnt.}

\subsection*{Ort \& Zeit}

\frage{2. Warum heißt der Ort \glqq Schädelstätte\grqq{}, und warum liegt er außerhalb der Stadt (V. 22)?}
\antwort{\glqq Golgatha\grqq\ ist aramäisch \emph{gulgalta} = Schädel; die Übersetzung steht im Text selbst (ebenso Mt 27,33). \textbf{Zur Namensherkunft gibt es drei Erklärungen:} (1) Am wahrscheinlichsten: Der Felsenhügel vor den Toren Jerusalems sah von weitem wie ein Schädel bzw. eine Glatze aus -- daher der Name als gewöhnlicher Ortsname, nicht als theologische Gründung. (2) Folkloristische Deutung späterer Jahrhunderte: Dort lägen die Schädel Hingerichteter -- unwahrscheinlich, denn jüdisches Reinheitsrecht (Berührung mit Knochen macht unrein) schloss aus, dass die Stadt an einem Knochenort vorbeiführte. (3) Nach Joh 19,41 lag dabei ein Garten mit einem neuen Grab; Schädel als Stein-/Gartenschmuck ist denkbar, bleibt aber Spekulation. \textbf{Wichtig ist der Textbefund:} Markus übersetzt den Namen -- die Leserschaft hört bei der ersten Nennung ein Sterbenssymbol, ohne dass der Text Symbolik erklärt. \textbf{Zur Lage außerhalb der Stadt:} Hinrichtungen lagen vor den Toren -- aus drei zusammenwirkenden Gründen: (a) \textbf{Jüdisches Reinheitsrecht} (der entscheidende Grund in Jerusalem): Berührung mit einer Leiche macht unrein (Num 19,11--16), ein Aufgehängter durfte nicht über Nacht bleiben (Dtn 21,22f); die Tora brachte Delinquenten deshalb \glqq hinaus vor das Lager\grqq\  (Num 15,35; Lev 24,14) -- Jerusalem lebte dies als \glqq hinaus vor das Tor\grqq{}. (b) \textbf{Römische Praxis}: Kreuzigungen standen an Straßen und Stadtgrenzen -- zur Abschreckung (tagelang sichtbar für alle Ein- und Ausreisenden; Josephus berichtet für das Jahr 70 von hunderten Kreuzigungen täglich in Sichtweite der Stadt), weil die Strafe als Inszenierung Publikum brauchte (so auch V. 29: \glqq die vorübergingen\grqq{}), und aus praktischen Gründen (Schreie, Vögel, Verwesung gehörten nicht in die Stadt). (c) \textbf{Soziale Entwertung}: Das Hinausführen markierte den sozialen Tod -- der Hingerichtete wurde wie ein Verfluchter aus der Gemeinschaft entfernt. Hebr 13,12f nimmt das auf und kehrt es um: \glqq Jesus hat außer dem Tor gelitten, um das Volk zu heiligen -- so lass auch uns zu ihm hinausgehen aus dem Lager\grqq{}. Für Markus 15 heißt das: Das \glqq Hinausführen\grqq\  (V. 20.22) ist Teil der Strafe selbst -- und theologisch aufladbar: Jesus stirbt draußen, am Rand, im Unreinen. Der Ort war zudem öffentlich zugänglich: Passanten können die Inschrift lesen, spotten, stehen bleiben. Für die Gruppe reicht: aramäisch \glqq Schädel\grqq{}, Name wohl nach der Hügelform, nicht nach echten Schädeln.}

\frage{3. Sammle die Zeitangaben (V. 25. 33. 34) -- wie verläuft die Zeitleiste? Welches Fest ist im Hintergrund?}
\antwort{Dritte Stunde (= 9 Uhr): Kreuzigung (V. 25). Sechste bis neunte Stunde (= 12--15 Uhr): Finsternis (V. 33). Neunte Stunde (= 15 Uhr): Todesruf und Tod (V. 34--37). Die Finsternis umschließt exakt die Mittagszeit. Hintergrund ist das Passahfest: Jerusalem ist voller Pilger, die Stimmung politisch aufgeladen -- daher die Barabbas-Szene (15,6--15) kurz zuvor.}

\subsection*{Kontext}

\frage{4. Was geschah unmittelbar vorher (Mk 14,43--15,19)?}
\antwort{Verhaftung in Gethsemane (14,43--52; Judas verrät mit Kuss, die Jünger fliehen), Verhör vor dem Hohen Rat mit Jesu Bekenntnis \glqq Ich bin's\grqq\ und dem Urteil \glqq schuldig des Todes\grqq\ (14,53--65), Verleugnung Petris (14,66--72), Verhandlung vor Pilatus samt Barabbas-Ausschreibung (15,1--15), Geißelung und Verspottung durch die Soldaten mit Purpurmantel und Dornenkrone (15,16--19). Vers 20 schließt diese Szene ab: \glqq als sie ihn verspottet hatten \ldots\grqq\ }

\frage{5. \glqq Sohn Gottes\grqq{}: Mk 1,1 $\leftrightarrow$ V. 39 -- wer erkennt es am Ende, und warum ist das überraschend?}
\antwort{Mk 1,1 titelt Jesus als \glqq Sohn Gottes\grqq{}; danach bezeugen es Gott (Taufe 1,11; Verklärung 9,7) und Dämonen -- aber kein Mensch versteht es. Erst in 15,39 bekennt ein \textbf{Mensch} diesen Titel: ein römischer Offizier, Heide, selbst an der Hinrichtung beteiligt -- und ausgerechnet beim Anblick des Sterbenden, nicht eines Wunders. Gottessohnschaft offenbart sich bei Markus nicht in Macht, sondern am Kreuz.}

\subsection*{Figuren}

\frage{6. Wer handelt, wer redet, wer schweigt in V. 20--41? (Wann spricht Jesus selbst?)}
\antwort{Handelnd: die Soldaten (kreuzigen, verlosen die Kleider), Simon von Kyrene (trägt gezwungen das Kreuz), die Vorübergehenden, Hohenpriester und Schriftgelehrten und die Mitgekreuzigten (alle lästern), ein Umstehender mit dem Essigschwamm, der römische Hauptmann (bekennt), die Frauen (sehen von ferne). \textbf{Jesus redet nur ein einziges Mal selbst}: der Psalm-Ruf in V. 34. Im ganzen Abschnitt ist er Objekt -- ihm geschieht alles, er wirkt kaum. Das ist Erzählabsicht: Der Leidende, nicht der Kämpfende, steht im Zentrum.}

\frage{7. Wer fehlt -- und warum? (Jünger, Mk 14,50!)}
\antwort{Die Jünger. Sie sind bei der Verhaftung geflohen (14,50), Petrus hat verleugnet (14,66--72). Am Kreuz sind nur Gegner, Unbeteiligte und Außenstehende. Die Abwesenheit ist Teil der Erzählung: Am Ende erkennen ausgerechnet Fremde (Zenturio, Frauen), was die Vertrauten nicht erkannt haben.}

\section*{Kreuzigung (V. 20--28)}

\frage{8. Simon wird \textbf{gezwungen}, das Kreuz zu tragen (V. 21) -- was heißt das für Mk 8,34 (\glqq sein Kreuz aufnehmen\grqq{})?}
\antwort{Simon trägt das Kreuz unfreiwillig, gepresst -- das genaue Gegenteil des Nachfolge-Wortes von Mk 8,34, wo das Kreuz \emph{freiwillig} \glqq auf sich genommen\grqq\ wird. Simon ist der Prototyp des unfreiwilligen Kreuztragens. Und doch: Er geht den Weg \emph{mit} Jesus. Dass sein Sohn Rufus der Gemeinde bekannt ist (Röm 16,13), deutet darauf, dass aus dem Gezwungenen ein Christ wurde. Nachfolge beginnt oft unfreiwillig -- aber sie beginnt.}

\frage{9. Warum weist Jesus den Myrrhenwein ab (V. 23)?}
\antwort{Der Myrrhenwein ist ein Betäubungsmittel gegen die Schmerzen (vgl. Spr 31,6), wohl eine Geste der Barmherzigkeit frommer Jerusalemer Frauen. Jesus lehnt ab: Er will die Qual bewusst durchleiden -- auch den Schrei der Gottesverlassenheit (V. 34). Wer betäubt ist, kann nicht klagen. Der Leidensweg soll nicht betäubt, sondern durchlaufen werden.}

\frage{10. \glqq König der Juden\grqq\ (V. 26) -- als Anklage geschrieben: Inwiefern stimmt es?}
\antwort{Der titulus nannte das Delikt; \glqq König der Juden\grqq\ ist als politische Anklage (Aufruhr gegen Rom) gemeint. Aber die Anklage sagt die Wahrheit: Jesus \emph{ist} der König -- nur anders, als Macht erwartet. Der Titel fällt in Kap. 15 sechsmal (15,2.9.12.18.26.32): erst von Pilatus, dann spöttisch, schließlich als Inschrift -- Markus macht aus der Verhöhnung eine Christologie. Der König thront am Kreuz.}

\section*{Verspottung am Kreuz (V. 29--32)}

\frage{11. \emph{Unser Leben}: Wo erleben wir, dass uns etwas zugemutet wird, das wir nicht gewählt haben -- und wie kann gerade darin Nachfolge beginnen?}
\antwort{Offene Frage. Mögliche Impulse: Krankheit, Arbeitslosigkeit, Pflege eines Angehörigen, Einschränkungen, die niemand wählt -- Simon trägt das Kreuz, ohne gefragt zu werden. Der Text entlastet: Nachfolge beginnt nicht mit Heldentum, sondern oft im Getragenwerden. Frage an die Gruppe: Gab es einen Moment, in dem aus einem Unfreiwilligen ein Freiwilliger wurde? (Rufus, Röm 16,13, deutet darauf.)}

\frage{12. Drei Gruppen spotten -- was wirft jede vor, wo wiederholt sich etwas?}
\antwort{Vorübergehende: Tempelzerstörung/Neubau in drei Tagen (V. 29f) -- Angriff auf sein Ansehen als Prophet. Hohenpriester und Schriftgelehrte: \glqq Er hat andern geholfen, sich selber nicht\grqq\ (V. 31) -- Angriff auf seine rettende Macht. Die Mitgekreuzigten: \glqq Der Christus, der König von Israel\grqq\ (V. 32) -- Angriff auf seinen Messiasanspruch. Sich wiederholt: Alle drei fordern ein sichtbares Machtzeichen (\glqq steig herab\grqq{}) als Beweis, und alle drei nennen Titel Jesu -- in Spottform.}

\frage{13. \textbf{Das Paradoxon von V. 31}: \glqq Er hat andern geholfen -- und kann sich selber nicht helfen.\grqq\ Wie kann beides zugleich wahr sein? Ist Nicht-sich-helfen-Können hier Schwäche -- oder gerade die Art, wie er anderen hilft?}
\antwort{Der Spott denkt in menschlicher Logik: Wer Macht hat, nutzt sie zuerst für sich -- sich selbst nicht helfen zu können beweist, dass die Hilfe für andere wertlos war. Das Kreuz dreht diese Logik um: Jesus hilft anderen \emph{dadurch}, dass er sich nicht selbst hilft -- die Lebenshingabe ist die Hilfe (Mk 10,45: \glqq sein Leben zu geben als Lösegeld für viele\grqq{}). Das Paradox ist keine Logiklücke, sondern der Kern der Botschaft: Rettung geschieht nicht \emph{trotz} des Kreuzes, sondern \emph{durch} es. Die Hohenpriester sprechen unfreiwillig aus, was sie nicht meinen: Genau daran, dass er sich nicht selbst hilft, erkennt der Hauptmann wenig später den Gottessohn (V. 39). Für heute: Auch christliche Hilfe geschieht oft nicht aus Übermacht, sondern aus Hingabe -- und wird genau deshalb von Weltlogik als Schwäche verwechselt.}

\frage{14. Wo wird der Spott unfreiwillig wahr (V. 29. 31. 32)?}
\antwort{In allen drei Fällen -- nur anders, als die Spötter meinen: (a) Den Tempel \glqq bricht er ab und baut auf in drei Tagen\grqq{}: sein Tod und seine Auferweckung (so gedeutet in Joh 2,19--21). (b) Er \glqq hilft andern\grqq{}: gerade dadurch, dass er sich nicht selbst hilft -- das Kreuz wird zur Rettung der Vielen. (c) Er \emph{ist} Christus und König Israels -- als Gekreuzigter. Markus lässt die Gegner unfreiwillig die Wahrheit verkündigen: Die Leser erkennen, was die Spötter nicht ahnen.}

\frage{15. \glqq \ldots\ damit wir sehen und glauben\grqq\ (V. 32) -- warum ist das keine echte Glaubensbedingung?}
\antwort{Weil Glaube, der ein sichtbares Wunder verlangt, kein Vertrauen mehr ist, sondern Wissenszwang. Das Kreuz verweigert das Zeichen bewusst: Der Glaube, den Markus sucht, entsteht am gekreuzigten Jesus selbst -- ohne Absicherung (vgl. Joh 20,29: \glqq Selig sind, die nicht sehen und doch glauben\grqq{}). Erst nach Ostern wird das Kreuz als Zeichen lesbar -- aber es wird nie zum Beweis.}

\section*{Finsternis, Schrei und Bekenntnis (V. 33--41)}

\frage{16. \emph{Unser Leben}: Wo erwarten wir (von uns selbst oder anderen) erst Beweise -- und was hieße glauben ohne Beweis?}
\antwort{Offene Frage. Mögliche Impulse: Beziehungen leben wir ohne Beweise -- Vertrauen wächst in Handlungen, nicht in Beweisführungen. Wo verlangen wir von Gott Zeichen (Erhörung, Wendung), bevor wir vertrauen? Das Kreuz verweigert den Beweis und schenkt trotzdem Gemeinschaft. In der Gruppe behutsam bleiben: nicht schuldig sprechen, sondern ehrlich über das Bedürfnis nach Sicherheit reden.}

\frage{17. \glqq Warum hast du mich verlassen?\grqq\ (V. 34) -- nur Verzweiflung? (Ps 22: Klage, die ins Lob mündet.)}
\antwort{Beides. Erstens echte Gottesverlassenheit: Jesus durchleidet die Distanz zum Vater -- im Neuen Testament gedeutet als Tragen der Sünde (2 Kor 5,21; Gal 3,13). Zweitens ein Glaubenswort: Der Ruf zitiert Ps 22,2 -- und dieser Psalm mündet ins Lob und in die Verkündigung unter den Völkern (Ps 22,23--32). Wer den Psalm kennt, hört im Schrei auch Vertrauen: Jesus redet Gott trotz allem mit \glqq meinem Gott\grqq\ an. Klage und Glaube schließen sich hier nicht aus.}

\frage{18. Finsternis und Vorhangriss (V. 33. 38) -- wer handelt hier, was wird geöffnet bzw. gerichtet?}
\antwort{Gott selbst -- beide Zeichen geschehen ohne Zutun Jesu. Die Finsternis über dem Land (12--15 Uhr) ist Gerichts- und Theophanie-Zeichen (Am 8,9; Ex 10,21f). Der Vorhang vor dem Allerheiligsten reißt \glqq von oben bis unten\grqq\ -- von Gott her. Doppelt zu deuten: (a) Der Zugang zu Gott wird für alle geöffnet (Hebr 10,19f: \glqq durch den Vorhang\grqq{}); (b) Gericht über den Tempelkult -- seine Zeit ist abgelaufen (vgl. Mk 13,2). Beides zusammen: Die alte Trennung endet, Gott beginnt woanders.}

\frage{19. Warum bekennt ausgerechnet der römische Hauptmann \glqq Gottes Sohn\grqq\ (V. 39)?}
\antwort{Er ist Heide, Offizier der Besatzungsmacht, verantwortlich für die Vollstreckung -- der letzte Mensch, von dem man es erwarten würde. Sein Anlass ist nicht ein Wunder, sondern \emph{wie} Jesus stirbt: \glqq daß er mit solchem Geschrei verschieden war\grqq{}. Der Sterbende selbst überzeugt ihn. Damit erfüllt sich Mk 1,1: Am Kreuz erkennt ein Mensch, wer Jesus ist. Dass es ein Feind tut, macht die Pointe scharf.}

\frage{20. Warum bleiben die Frauen und sehen alles (V. 40f) -- anders als die Jünger?}
\antwort{Die Frauen stehen \glqq von ferne\grqq\ -- aber sie \emph{bleiben}, während die Jünger flohen (14,50). Sie waren schon in Galiläa nachgefolgt und haben ihm gedient (V. 41) -- Markus nennt sie damit in Jünger-Terminologie. Ihr Bleiben hat einen praktisch-theologischen Grund: Sie sind die einzigen Zeuginnen von Tod \emph{und} Begräbnis (V. 47) -- und deshalb die glaubwürdigen Zeuginnen des leeren Grabes (16,1--8). Die Osterbotschaft des Evangeliums steht auf ihrem Zeugnis.}

\frage{21. \emph{Unser Leben}: Wann haben wir Gottesnähe oder -verlassenheit am eigenen Leib erlebt? Was half in der Finsternis?}
\antwort{Offene Frage -- hier ist Raum für persönliche Geschichten. Wichtig: Klage ist biblisch (Ps 22, Hiob, Jeremia) und kein Glaubensmangel. Was half: Menschen, die \glqq von ferne\grqq\ blieben; Rituale (Gebet, Lied, Abendmahl); das Wissen um den Ps-22-Bogen (Klage endet nicht im Abgrund). Wer keine eigene Erfahrung teilen will, darf schweigen -- die Frage lädt ein, zwingt nicht.}

\frage{22. \emph{Unser Leben}: Für wen können wir heute \glqq von ferne\grqq\ da sein -- dableiben, zusehen, Zeugnis geben?}
\antwort{Offene Frage. Mögliche Impulse: Menschen im Krankenhaus, in der Isolation, im Sterben; Nachbarn in Einsamkeit; auch hier im Hauskreis: Wer ist gerade \glqq von ferne\grqq\ (verletzt, ausgezogen, verschlossen)? Die Frauen beweisen: Distanz ist kein Hindernis für Zeugenschaft -- da sein reicht. Überleitung zu den Gebetsanliegen (19 Uhr).}

\section*{Querbezüge zu anderen Bibelstellen}

\frage{23. Lies Psalm 22 -- welche Verse unseres Textes klingen dort an?}
\antwort{Ps 22,2 → V. 34 (\glqq Mein Gott, mein Gott, warum hast du mich verlassen?\grqq{}); Ps 22,8 → V. 29--32 (\glqq Er vertraute auf den HERRN, der möge ihn erretten\grqq\ -- so klingt der Spott der Vorübergehenden); Ps 22,19 → V. 24 (\glqq Sie verteilen unter sich meine Kleider und werfen das Los um mein Gewand\grqq{}). Der Passionsbericht ist von Anfang an als Auslegung dieses Psalms erzählt: Das Leiden als Gebet, die Klage als Glauben.}

\frage{24. Jesaja 53: Welche Züge des \glqq leidenden Gottesknechtes\grqq\ (Jes 53,5.12) wiederholen sich in unserem Abschnitt?}
\antwort{Jes 53,5 (\glqq um unserer Übertretungen willen zerschlagen\grqq{}) -- das stellvertretende Leiden, das der Schrei und der Tod andeuten. Jes 53,12 (\glqq er ist unter die Übeltäter gerechnet\grqq{}) -- wörtlich erfüllt in V. 27f, zwei Räuber rechts und links. Jes 53,7 (stillschweigend wie ein Lamm) -- der schweigende Jesus von V. 20 bis 37, der nur ein Wort spricht.}

\frage{25. Mk 8,31; 9,31; 10,33f: Jesus kündigt sein Leiden dreimal an -- was davon wird in V. 20--27 buchstäblich erfüllt?}
\antwort{Die Ankündigungen enthalten die Sequenz: übergeben -- verurteilt -- den Heiden ausgeliefert -- verspottet -- ausgepeitscht -- getötet. Genau diese Sequenz läuft in 15,1--27 in derselben Wortwahl ab (\glqq ausliefern\grqq{}, \glqq verurteilen\grqq{}, \glqq verspotten\grqq{}, \glqq geißeln\grqq{}, \glqq kreuzigen\grqq{}). Markus zeigt: Das Kreuz ist kein tragisches Unglück, sondern geschieht nach Jesu eigenem Wegverständnis -- dreimal angekündigt, \glqq nach der Schrift\grqq\ erfüllt.}

\frage{26. Mk 10,45: \glqq \ldots\ sein Leben zu geben als Lösegeld für viele\grqq\ -- wie hört sich dieses Wort im Licht von V. 37.39 an?}
\antwort{In V. 37 gibt Jesus sein Leben: \glqq Jesus aber rief laut und verschied.\grqq\ Der Rahmen lautet: Lebenshingabe (V. 37) -- Bekenntnis (V. 39). Was die Lebenshingabe bewirkt, zeigt die erste Frucht: Ausgerechnet der Feind kommt zum Glauben. Das Wort von Mk 10,45 wird am Kreuz erzählerisch eingelöst -- als Deutewort, nicht als Rechenformel.}

\frage{27. Hebr 10,19f: Was bedeutet der Vorhangriss (V. 38) im Licht dieses neutestamentlichen Deuteworts?}
\antwort{Hebr 10,19f: \glqq Wir haben Zuversicht, einzugehen in das Heiligtum durch das Blut Jesu, den neuen und lebendigen Weg, den er uns geweiht hat durch den Vorhang.\grqq\ Der Vorhang ist dort Bild für Jesu Fleisch/Tod: Sein Tod öffnet den Weg zu Gott. Der Riss \glqq von oben bis unten\grqq\ (V. 38) sagt dieselbe Sache erzählerisch: Der Zugang steht offen -- von Gott her geöffnet, für jeden zugänglich, nicht mehr durch den Kult, sondern durch Christus.}

\frage{28. Amos 8,9 / Ex 10,21f: Welches alte Motiv steht hinter der Finsternis (V. 33)?}
\antwort{Am 8,9: \glqq An jenem Tag lasse ich die Sonne zur Mittagszeit untergehen und lasse am hellen Tage das Land finster werden\grqq\ -- Gerichtsansage für den Tag des Herrn. Ex 10,21f: die Finsternis als neunte Plage über Ägypten -- Gerichtszeichen. Beide zusammen: Die Finsternis von 12--15 Uhr signalisiert, dass hier Gottes Tag anbricht -- Gericht und Heilsentscheidung fallen in dieselben drei Stunden.}

\section*{Abschluss}

\frage{29. Der Hauptmann erkennt Jesus im Sterbenden. Wo begegnen uns heute Menschen in \glqq gekreuzigten\grqq\ Situationen?}
\antwort{Offene Frage -- mögliche Impulse für die Gruppe: Menschen in Krankheit, Einsamkeit, Schuld, Verzweiflung; Sterbende, die niemand mehr besucht; Menschen am Rand. Der Blick des Hauptmanns kann Maßstab werden: \emph{hinsehen} statt wegschauen, \emph{erkennen} statt urteilen. Das leitet direkt in den Gebetsanliegen-Teil über (19 Uhr, 2 Gruppen).}

\vspace{1em}
{\footnotesize Leitfaden für die Leitung · Hauskreis 08.10.2026 · Bibeltext: Lutherbibel 2017, © Deutsche Bibelgesellschaft.}
\end{document}
